% Part: incompleteness
% Chapter: representability-in-q
% Section: introduction

\documentclass[../../../include/open-logic-section]{subfiles}

\begin{document}

\olfileid{inc}{req}{int}
\olsection{પરિચય}

અપૂર્ણતાનાં પ્રમેયો એવા સિદ્ધાંતોને લાગુ
પડે છે જેમાં સંગણનીય વિધેયો વિશેનાં
પાયાનાં તથ્યો વ્યક્ત અને સાબિત કરી
શકાય. આવા અત્યંત ન્યૂનતમ સિદ્ધાંતનું
વર્ણન કરીશું; તેનું નામ
``$\Th{Q}$'' છે (રાફેલ રોબિન્સનના
નામ પરથી તેને ક્યારેક
``રોબિન્સનનું $Q$'' પણ કહે છે).
વિધેય $\Th{Q}$માં \emph{નિરૂપ્ય}
હોય તેનો અર્થ જણાવીશું અને પછી આ
સાબિત કરીશું:
\begin{quote}
  વિધેય $\Th{Q}$માં નિરૂપ્ય હોય
  ત્યારે અને માત્ર ત્યારે તે સંગણનીય છે.
\end{quote}
એક રીતે આ આપણને સંગણનીયતાનું બીજું
નિદર્શ આપે છે. પરંતુ અટકવાની
સમસ્યાને તેમાં રૂપાંતરિત કરીને
$\Setabs{!A}{\Th{Q} \Proves !A}$
ગણ નિર્ણેય નથી તે બતાવવા પણ તેને
વાપરીશું. અંતે આના કરતાં ઘણાં વધુ
શક્તિશાળી પરિણામો સાબિત કર્યાં હશે.

$\Th{Q}$ની ભાષા અંકગણિતની ભાષા છે.
$\Th{Q}$ નીચેના સ્વયંસિદ્ધ વાક્યોનું
બનેલું છે (તેમની સાથે
!!{identity}વાળા પ્રથમ-ક્રમના
તર્કશાસ્ત્રનાં બીજા સ્વયંસિદ્ધ વાક્યો
અને નિયમો પણ વાપરવાના છે):
\begin{align*}
& \lforall[x][\lforall[y][(\eq[x'][y'] \lif \eq[x][y])]] \tag{$!Q_1$}\\
& \lforall[x][\eq/[\Obj 0][x']] \tag{$!Q_2$}\\
& \lforall[x][(\eq[x][\Obj 0] \lor \lexists[y][\eq[x][y']])] \tag{$!Q_3$}\\
& \lforall[x][\eq[(x + \Obj 0)][x]] \tag{$!Q_4$}\\
& \lforall[x][\lforall[y][\eq[(x + y')][(x + y)']]] \tag{$!Q_5$}\\
& \lforall[x][\eq[(x \times \Obj 0)][\Obj 0]] \tag{$!Q_6$}\\
& \lforall[x][\lforall[y][\eq[(x \times y')][((x \times y) + x)]]] \tag{$!Q_7$}\\
& \lforall[x][\lforall[y][(x < y \liff \lexists[z][\eq[(z' + x)][y]])]] \tag{$!Q_8$}
\end{align*}
દરેક પ્રાકૃતિક સંખ્યા~$n$ માટે
અંકનિરૂપણ~$\num{n}$ને એવું
પદ~$\Obj{0}^{\prime\prime\ldots\prime}$
વ્યાખ્યાયિત કરીએ જેમાં કુલ~$n$
ઉપરચિહ્નો હોય. આમ $\num{0}$
માત્ર !!{constant}~$\Obj{0}$ છે,
$\num{1}$ એ $\Obj{0}'$ છે,
$\num{2}$ એ $\Obj{0}''$ છે, વગેરે.

અંકગણિતના સિદ્ધાંત તરીકે
$\Th{Q}$ \emph{અત્યંત} નબળો છે.
દાખલા તરીકે,
$\lforall[x][\eq/[x][x']]$ અથવા
$\lforall[x][\lforall[y][(x + y) = (y +
    x)]]$ જેવાં ઘણાં સરળ વિધાનો પણ
તેમાં સાબિત થઈ શકતાં નથી. જોકે
$\Th{Q}$ આટલો નબળો છે \emph{એટલા
માટે જ} તે રસપ્રદ છે, એવું આપણે
જોઈશું. હકીકતમાં અપૂર્ણતા પ્રમેય
લાગુ પડે એટલો તે માંડ શક્તિશાળી છે.
$\Th{Q}$ રસપ્રદ હોવાનું બીજું કારણ
તેનો સ્વયંસિદ્ધ વાક્યોનો ગણ
\emph{સાન્ત} હોવો છે.

$\Th{Q}$માં આગમનની યોજના ઉમેરીએ,
તો $\Th{Q}$ કરતાં વધુ શક્તિશાળી સિદ્ધાંત
(\emph{પિયાનો અંકગણિત}~$\Th{PA}$)
મળે છે:
\[
(!A(\Obj 0) \land \lforall[x][(!A(x) \lif !A(x'))]) \lif \lforall[x][!A(x)]
\]
અહીં $!A(x)$ કોઈપણ સૂત્ર હોઈ શકે.
$!A(x)$માં $x$ સિવાયના મુક્ત
!!{variable}s હોય, તો એ બધાંને બાંધવા
આગળ સર્વપરિમાણકો ઉમેરીએ છીએ
(જેથી આગમન-યોજનાનું અનુરૂપ દૃષ્ટાંત
!!a{sentence} બને). દાખલા તરીકે,
$!A(x, y)$માં !!{variable}~$y$
પણ મુક્ત હોય, તો અનુરૂપ દૃષ્ટાંત આ છે:
\[
\lforall[y][((!A(\Obj 0) \land \lforall[x][(!A(x) \lif !A(x'))]) \lif
  \lforall[x][!A(x)])]
\]
આગમન-યોજનાનાં દૃષ્ટાંતો વાપરીને
$\Th{PA}$નાં સ્વયંસિદ્ધ વાક્યોમાંથી
$\Th{Q}$નાં વાક્યો કરતાં ઘણું વધુ
સાબિત કરી શકાય છે. હકીકતમાં,
પ્રાકૃતિક સંખ્યાઓ વિશેનાં એવાં
``સ્વાભાવિક'' વિધાનો શોધવા ખાસ્સી
મહેનત પડે છે જે પિયાનો અંકગણિતમાં
સાબિત ન થઈ શકે!{}

\begin{defn}
\ollabel{defn:representable-fn}
  પ્રાકૃતિક સંખ્યાઓમાંથી પ્રાકૃતિક
  સંખ્યાઓમાં જતું વિધેય
  $f(x_0,\ldots,x_k)$
  {$\Th{Q}$માં \emph{નિરૂપ્ય}}
  કહેવાય, જો એવું સૂત્ર
  $!A_f(x_0,\dots,x_k,y)$ હોય કે
  જ્યારે પણ $f(n_0,\dots,n_k) = m$
  હોય, ત્યારે $\Th{Q}$ નીચેના
  બંને વિધાનો સાબિત કરે:
\begin{enumerate}
\item\ollabel{defn:rep:a} $!A_f(\num{n_0}, \dots, \num{n_k}, \num{m})$
\item\ollabel{defn:rep:b} $\lforall[y][(!A_f(\num{n_0}, \dots,
\num{n_k}, y) \lif \num{m} = y)]$.
\end{enumerate}
\end{defn}

વ્યાખ્યા બીજી રીતે પણ કહી શકાય.
દાખલા તરીકે, સમકક્ષ રીતે એવી
શરત મૂકી શકાય કે $\Th{Q}$
$\lforall[y][(!A_f(\num{n_0}, \dots,
\num{n_k}, y) \liff \eq[y][\num{m}])]$
સાબિત કરે.

\begin{thm}
\ollabel{thm:representable-iff-comp}
વિધેય $\Th{Q}$માં નિરૂપ્ય હોય ત્યારે
અને માત્ર ત્યારે તે સંગણનીય છે.
\end{thm}

પ્રમેય સાબિત કરવાની બે દિશાઓ છે.
વાક્યરચનાનું અંકગણિતીકરણ હાથવગું
હોય, તો ડાબેથી જમણી દિશા ઘણી
સીધી છે. બીજી દિશા વધુ મહેનત માગે
છે. મૂળ વિચાર આ છે: ``સંગણનીય''નો
ચોક્કસ અર્થ આપવા ``સામાન્ય
પુનરાવર્તી'' પસંદ કરીએ અને દરેક
સામાન્ય પુનરાવર્તી વિધેય
$\Th{Q}$માં નિરૂપ્ય છે તે બતાવીએ.
યાદ કરો કે વિધેય સામાન્ય પુનરાવર્તી
છે જો તે $\Zero$, ઉત્તરગામી
વિધેય~$\Succ$ અને પ્રક્ષેપ
વિધેયો~$\Proj{n}{i}$માંથી સંયોજન,
આદિમ પુનરાવર્તન અને નિયમિત
લઘુત્તમીકરણ વાપરીને વ્યાખ્યાયિત
કરી શકાય. તેથી દરેક સામાન્ય
પુનરાવર્તી વિધેય $\Th{Q}$માં
નિરૂપ્ય છે તે બતાવવાની એક રીત
એ છે કે પાયાનાં વિધેયો નિરૂપ્ય છે
તે સાબિત કરીએ અને નિરૂપ્ય
વિધેયોમાંથી સંયોજન, આદિમ
પુનરાવર્તન તથા નિયમિત
લઘુત્તમીકરણ વડે વ્યાખ્યાયિત
થતાં વિધેયો પણ નિરૂપ્ય છે તે
સાબિત કરીએ. બીજા શબ્દોમાં,
પાયાનાં વિધેયો નિરૂપ્ય છે અને
નિરૂપ્ય વિધેયોનો વર્ગ સંયોજન,
આદિમ પુનરાવર્તન તથા નિયમિત
લઘુત્તમીકરણ હેઠળ ``બંધ'' છે
તે બતાવી શકીએ. આમ દરેક સામાન્ય
પુનરાવર્તી વિધેયની નિરૂપ્યતા મળે.

પરંતુ નિરૂપ્ય વિધેયો આદિમ
પુનરાવર્તન હેઠળ બંધ છે તે બતાવવાનું
પગલું મુશ્કેલ છે. તેને ટાળવા
પહેલાં બતાવીએ કે ખરેખર આદિમ
પુનરાવર્તન વિના પણ કામ ચાલે છે.
એટલે કે દરેક સામાન્ય પુનરાવર્તી
વિધેય પાયાનાં વિધેયોમાંથી માત્ર
સંયોજન અને નિયમિત લઘુત્તમીકરણ
વાપરીને વ્યાખ્યાયિત થઈ શકે છે.
આ માટે બતાવીએ કે આદિમ પુનરાવર્તન
કોઈ ચોક્કસ નિયમિત લઘુત્તમીકરણ
વડે કરી શકાય છે. જોકે આ રીત
કાર્ય કરે તે માટે પાયાનાં વિધેયોમાં
કેટલાક વધુ વિધેયો ઉમેરવા પડે:
સરવાળો, ગુણાકાર અને સમતા
સંબંધનું લાક્ષણિક વિધેય~$\Char{=}$.
પછી આ \emph{બધાં} પાયાનાં વિધેયો
$\Th{Q}$માં નિરૂપ્ય છે અને નિરૂપ્ય
વિધેયો સંયોજન તથા નિયમિત
લઘુત્તમીકરણ હેઠળ બંધ છે તે
બતાવીને પ્રમેય સાબિત કરી શકીએ.

\end{document}
